\thispagestyle{plain}
\begingroup\centering
{\fontsize{9.5}{10.7}\selectfont Manuscrit de compilation pour le transfert académique\par}
\vspace{3mm}
{\fontsize{17}{20.5}\selectfont\bfseries\Titulo\par}
\vspace{6pt}
{\fontsize{11.5}{14}\selectfont\Autores\par}
\vspace{3pt}
{\fontsize{8}{10}\selectfont Coautorat sans préséance de contribution\par}
\vspace{3pt}
{\fontsize{8}{10}\selectfont Intégration et compilation documentaires générées par ChatGPT - Astra.\par}
\par\endgroup
\vspace{5pt}
{\centering\bfseries Résumé\par}
\vspace{2pt}
\begingroup
\fontsize{9.5}{10.7}\selectfont
\setlength{\parskip}{1.5pt plus .5pt minus .5pt}
\noindent
On génère coinductivement \(\pi,\varphi,e\) conjointement et à profondeur arbitraire, on
dérive la constante de structure fine \(\alpha\) et on construit la
coordonnée électronique, sa représentation de spin \(1/2\) et ses
composantes d'hélicité. La construction relie la détermination de
ces valeurs à la conservation des structures qui les produisent.
Son fondement est l'\emph{holographie modulaire triadique} (HMT) :
\emph{All Possible Paths} (APP) est le support toroïdal discret
\((\mathbb Z/9\mathbb Z)^2\), muni du graphe de Cayley additif
de générateurs \(\pm(1,0),\pm(0,1)\) et des évaluations somme et
produit ; TRIT est la signature ternaire orientée \(\{-1,0,+1\}\), qui
distingue trois régimes quadratiques ; le \emph{Topological Phase Kernel}
(TPK) compose la sélection de phase, le transport et la mise à jour de la mémoire.

Les 104\,976 conditions initiales orientées produisent 468 émissions et 243
séquences ternaires visibles. Dans ce domaine sont déterminées des régions de
clôture, de propagation et d'auto-échelle, dont les réalisations archimédiennes sont
\(\pi,e,\varphi\). Le raffinement par cylindres compatibles permet d'obtenir
tout préfixe fini par un calcul fini. L'holonomie nonadique
conserve la phase tout en accroissant la mémoire ; sa représentation apériodique
explicite la coexistence d'une périodicité modulaire et d'un prolongement illimité.

Le registre intégral dodécaphasique \(K\) et sa réalisation rationnelle périodique
\(\kappa\) admettent des reconstructions réversibles. Leur composition avec les registres
régionaux détermine \(\alpha\) par normalisation positionnelle. La voie analytique
corrélée la représente comme unique racine simple d'un prolongement
de la résolvante des vacances, avec une récurrence et une classe géométrique explicites.
La compatibilité à toute profondeur démontre l'égalité des deux
représentations du même état. L'évaluation scalaire et la résolution d'incidence
s'articulent par les registres orientés que toutes deux conservent.

La construction électronique réunit le centre d'APP, la vacance orientée,
une représentation matricielle des trois régimes et le transfert entre
les évaluations aréale et volumique. À partir du bloc électronique sont construits
les générateurs de spin \(1/2\) et les projecteurs directionnels ; leur interprétation
comme hélicité déclare la direction de propagation et conserve l'évaluation
scalaire, l'orientation de charge et la mémoire de la route comme données distinctes.
Deux lecteurs du retour déterminent
les coefficients de sa composition complète. La trajectoire centrale visite
exactement \(27\) cellules ; ses historiques orientés se distinguent au moyen
de la coordinatisation intégrale réversible. La décomposition harmonique du retour
produit les modes \(2,6,10\), dont les puissances \(64,16,64\) restituent
les coefficients \(80\) et \(6\) du correcteur électronique.
L'échelle d'action et la
réalisation dimensionnelle sont exposées avant la comparaison avec CODATA 2018 et
2022. La seconde réalisation de la structure conserve l'incidence de
frontière : plans de Witt, codes de Golay, réseau de Leech et algèbre
de sommets \(V^\natural\), avec son groupe d'automorphismes, le Monstre.
Les réalisations numérique et exceptionnelle procèdent de la même structure ;
leurs opérations locales et leurs limites compatibles conservent les registres
qui relient équations, valeurs et incidence.
Le registre \(K\) sélectionne en outre la direction de la réduction
\(11\to10\). Les orbifolds d'ordres \(2\) et \(3\) fournissent
des présentations isomorphes de \(V^\natural\), avec transport de leurs
automorphismes et de leurs traces graduées. L'isomorphisme des écrans
quotientés entrelace la dualité des rayons réciproques, en conservant
les données de réduction de \(K\) et l'écran réticulaire.

% BEGIN S27_FRAMING
La comparaison des étoiles de \(B_K\), \(B_\varphi\) et
\(B_{\rm APP}\) distingue leurs douze complétions au moyen d'une polarisation
commune et situe leurs signatures dans la classification des \(495\) quadruplets.
Le cadre ordonné \((H^-_\alpha,H_e,H_\varphi,H_{\rm APP})\) possède un
stabilisateur trivial dans le groupe d'incidence construit. Le cocycle
de retenue explicite, à son tour, l'information des quotients entiers
qui complète les résidus des colonnes régionales. Ses lois de
composition et de transport conservent le relèvement arithmétique associé
à la lecture numérique, tandis que les étoiles et leurs stabilisateurs
déterminent l'organisation de la lecture de frontière.

% END S27_FRAMING
\vspace{3pt}
La monodromie du transport et le modèle mathématique de cristal temporel
apériodique articulent le retour modulaire, l'avancée de la mémoire et les trois
régimes de l'exponentielle. La décade déterminantale et le caractère
d'action déterminent la section correspondant à la constante de Planck
réduite, \(\hbar\), qui intervient dans la composition électronique.



% BEGIN R27_FRAMING_SUMMARY_ES
La précision de cette chaîne se développe à partir de l'arithmétique de l'émetteur :
on détermine son alphabet décimal de \(42\) blocs, on présente les raccordements
entre émissions régionales et on distingue la mémoire ordonnée de la distribution
marginale. La sélection de frontière distingue saturation, neutralité duale
et orientation. Les lois de raffinement APP, le sous-réseau d'indice
deux, l'addition nonadique avec retenue et la corrélation quadratique explicitent
les opérations intermédiaires. Dans la réalisation de niveau trois, on obtient,
par une identité de séries formelles, une divisibilité uniforme
optimale par \(3^9\).
% END R27_FRAMING_SUMMARY_ES

\noindent\textbf{Mots-clés :} dérivation structurelle des constantes fondamentales ;
Holographie Modulaire Triadique ; génération coinductive ; constante de structure fine ;
structure algébrique de l'électron ; holonomie ; incidence de frontière ;
constante de Planck réduite ; monodromie ; cristal temporel apériodique ;
symétries exceptionnelles ; Moonshine.
\par\endgroup
